\scanpage{239}
qui s'écrit aussi
\[
\diagram{\begin{tikzcd}[column sep=large,row sep=large]
H^i(\mathrm{Fil}^p C_{\mathrm{DR}}^*(X;S,t,\gamma^*))\arrow[r,"\sim"]\arrow[d]&H^i(X,\mathrm{Fil}^pS)\arrow[d]\\
H^i(\mathrm{Fil}^q C_{\mathrm{DR}}^*(X;S,t,\unsure{\gamma'|\gamma^*}))\arrow[r]&H^i(X,\mathrm{Fil}^qS
\end{tikzcd}}
\]
\margin{$q\leq p$}
\margin{\figure{\begin{tikzpicture}\node[draw,ellipse] {\uuline{Iso si}};\end{tikzpicture}} \uuline{$i\leq p$}}
\margin{\uline{iso} si $i\leq q$}

\uline{NB} Si on ne suppose pas que l'épi
\[
k\{T\}\longrightarrow\mathrm{Gr}^*(S)
\]
est un iso, on ne peut rien dire. Prenons p.ex. le cas extrême de dégénérescence
\[
t=0,\quad\text{donc }S\xrightarrow{\sim}k
\]
on a alors
\[
\mathrm{Gr}^p\widehat A_*^{\unsure{0}}\simeq\Gamma_k^p(\Psi_*)\ .
\]
donc
\[
\mathrm{Gr}^p(\widehat A_*^{\unsure{0}})\simeq\unsure{\Pi_*|K_*}(k,p)
\]
qui n'est flasque pour aucun $p\geq0$ (si $k\ne0$).

Prenons aussi le cas de $(\mathbb Z_p,p)$, la pd-filtration définie par $p$ est cofinale à la filtration $p$-adique si $p\ne2$ (si $p=2$, c'est la filtration constante de valeur
\[
\mathrm{Fil}^{\unsure{r|p}}\mathbb Z_2=2\mathbb Z_2\text{ si }\unsure{r|p}\geq1,\quad\text{donc}\quad\mathrm{Gr}^*(S)\simeq\mathrm{Gr}^0(S)\simeq\mathbb F_2
\]
dans ce cas, i.e. c'est encore le cas de dégénérescence\ldots), quel est exactement l'idéal $S^{[i]}$ engendré par les $\dfrac{p^j}{j!}$ pour $j\geq i$ ?
\margin{\deleted{$\unsure{S^{[i]}=(\frac{p^n}{n!}\mathbb Z_p)}$}\deleted{\unsure{--- i.e. quelle}}}
\deleted{\unsure{i.e. quelle est}}
\margin{238}
